% Source: https://artofproblemsolving.com/community/c3h2368374_length_of_side_of_equilateral_triangle_lying_on_parabola
% Stable identity: aops-topic-2368374-q01
\textbf{Slopes of chords.} The chord through $\left(s,s^2\right)$ and $\left(t,t^2\right)$ has slope $\dfrac{t^2-s^2}{t-s}=s+t$. So the sides $PQ$, $QR$, $RP$ have slopes $p+q$, $q+r$, $r+p$.

\textbf{Directions.} $PQ$ has slope $\sqrt3$, so it makes an angle of $60^\circ$ with the $x$-axis. In an equilateral triangle the other two sides make angles of $60^\circ$ with $PQ$, on either side, so (as lines) they make angles $0^\circ$ and $120^\circ$ with the $x$-axis, and they are different lines. Their slopes are $0$ and $-\sqrt3$, so $\{q+r,\ r+p\}=\{0,-\sqrt3\}$.

\textbf{Solving.} In either case, adding the three slopes gives $2(p+q+r)=\sqrt3+0-\sqrt3=0$, so $p+q+r=0$ and $r=-(p+q)=-\sqrt3$. Then $\{q,p\}=\{\sqrt3,0\}$, and the vertices are $(0,0)$, $\left(\sqrt3,3\right)$, $\left(-\sqrt3,3\right)$.

\begin{asy}
usepackage("amsmath");
size(5cm);
defaultpen(fontsize(10pt));
import graph;
real f(real x){return x^2;}
draw((-2.4,0)--(2.4,0),Arrow(6)); draw((0,-0.4)--(0,4.6),Arrow(6));
label("$x$",(2.4,0),E); label("$y$",(0,4.6),N);
draw(graph(f,-2.1,2.1),blue+linewidth(0.9));
pair P=(0,0), Q=(sqrt(3),3), R=(-sqrt(3),3);
draw(P--Q--R--cycle,red+linewidth(0.9));
dot(P); dot(Q); dot(R);
label("$P(0,0)$",P,SE); label("$Q(\sqrt3,3)$",Q,E); label("$R(-\sqrt3,3)$",R,W);
draw((P+Q)/2--(1.75,1.3),gray); label("slope $\sqrt3$",(1.75,1.3),E,red+fontsize(8pt));
draw((P+R)/2--(-1.75,1.3),gray); label("slope $-\sqrt3$",(-1.75,1.3),W,red+fontsize(8pt));
label("slope $0$",(0,3),N,red+fontsize(8pt));
\end{asy}

The side length is the distance from $\left(-\sqrt3,3\right)$ to $\left(\sqrt3,3\right)$:
\[QR=2\sqrt3,\qquad\text{and indeed}\quad PQ=\sqrt{3+9}=2\sqrt3.\]
So the side length is $\boxed{2\sqrt3}$.
